\section{Ergodic random fields}

The shift group acts by the deterministic kernels $\hat\theta_i(A \mid \omega) = \1_A(\theta_i\omega)$
of (14.4), so that the shift-invariant fields $\Prob_\Theta$ of (14.1) are the $\Pi$-invariant
measures of Corollary (7.4) and the almost surely invariant events $\mathcal{I}(\mu)$ of (14.3) are
$\mathcal{I}_\Pi(\mu)$. Everything below is stated for a measurable action of an arbitrary group
$G$ on $(\Cfg, \mathcal{F})$.

\begin{definition}[Invariant fields and invariant events, Georgii (14.1)--(14.4)]
    \label{def:erg-invariant}
    \uses{def:ext-ae-invariant}
    \lean{ProbabilityTheory.Kernel.smulKernel,
      ProbabilityTheory.Kernel.smulInvariantMeasure_iff_forall_invariant,
      ProbabilityTheory.Kernel.measurableSet_aeInvariantSigmaAlgebraFamily_smul}
    \leanok

    For a group $G$ acting measurably on $(\Cfg, \mathcal{F})$ and $g \in G$, the kernel of the
    action of $g$ is the deterministic kernel $\hat g(A \mid \omega) = \1_A(g\omega)$, Georgii's
    $\hat\theta_i$. A measure is invariant under the action if and only if it is invariant under
    every $\hat g$, and for a finite invariant $\mu$ the $\mu$-almost surely invariant
    $\sigma$-algebra $\mathcal{I}_\Pi(\mu)$ of the family $\Pi = \set{\hat g : g \in G}$ is
    Georgii's $\mathcal{I}(\mu)$:
    \[\mathcal{I}(\mu) = \set{A \in \mathcal{F} : g^{-1}A = A \ \mu\text{-a.e. for all } g \in G}.\]
\end{definition}

\begin{theorem}[Ergodicity is extremality, Georgii (14.5)(a) and (14.6)]
    \label{thm:erg-145a}
    \uses{def:erg-invariant, thm:ext-74}
    \lean{ProbabilityTheory.Kernel.ergodicSMul_iff_mem_extremePoints}
    \leanok

    Let a group $G$, subject to no countability assumption, act measurably on
    $(\Cfg, \mathcal{F})$ and let $\mu$ be an invariant probability measure. Then $\mu$ is
    \emph{ergodic} --- trivial on $\mathcal{I}(\mu)$, which is Georgii's Definition (14.6) --- if
    and only if $\mu$ is an extreme point of the convex set of invariant probability measures.
    This is Corollary (7.4) applied to the kernels of (14.4).
\end{theorem}

\begin{remark}[Parts (b), (c) and the ergodic decomposition]
    \label{rmk:erg-145bc}
    \uses{thm:erg-145a, thm:ext-invariant-density}

    (14.5)(b): for invariant $\mu$ and $\nu \ll \mu$, the measure $\nu$ is invariant if and only if
    $\nu = f\mu$ with $f$ measurable for $\mathcal{I}(\mu)$; this is
    Theorem~\ref{thm:ext-invariant-density} at the kernels of (14.4). (14.5)(c): each
    $\mu \in \Prob_\Theta$ is determined within $\Prob_\Theta$ by its restriction to
    $\mathcal{I}$. The shift-specific development --- (14.7) ergodicity as cube-averaged mixing,
    (14.10) the ergodic decomposition of $\Prob_\Theta$, (14.12) density of the extreme boundary
    --- instantiates the decomposition machinery of Chapter 7.
\end{remark}
